\section{不确定性引导的模型精炼}

\subsection{方法论定位}

标准的主动学习（Active Learning）通常需要获取新的\textbf{真实标签}
（如新CFD仿真结果或风洞实验数据）来迭代改进模型。
本课题受限于数据来源（公开数据集），
无法获取新的真实气动仿真标签。

因此，本模块实现的方法在严格意义上属于
\textbf{不确定性引导的数据增强（Uncertainty-Guided Data Augmentation）}
或\textbf{基于检索的模型精炼（Retrieval-based Refinement）}：
利用DeepEnsemble的$\sigma(x)$不确定性估计
识别代理模型的薄弱区域，
从已有训练集中基于相似性检索补充相关样本，
通过模型重训练降低这些区域的预测不确定性。
这一机制虽然没有引入新的物理标签，
但有效利用了已有数据中尚未被充分学习的分布信息。

\subsection{单轮反馈精炼机制}

在第5章构建的风险敏感PSO基础上，
进一步引入不确定性引导的单轮反馈精炼机制：
通过PSO探索反馈迭代精炼代理模型。
该机制的核心步骤为：

\begin{enumerate}
    \item \textbf{PSO探索}：风险敏感PSO在参数空间内搜索最优解，
        同时收集候选位置的$\sigma(x)$不确定性信息；
    \item \textbf{不确定性反馈}：DeepEnsemble估计每个访问位置的预测标准差，
        识别模型不可靠的高不确定性区域；
    \item \textbf{候选选择}：选取Top-K个$\sigma(x)$最高的候选位置，
        这些位置代表代理模型需要增强的"薄弱区"；
    \item \textbf{数据增强}：在高不确定性候选位置附近，
        从原始训练集中选取最相似的样本加入训练集；
    \item \textbf{模型精炼}：在扩充后的训练集上重新训练DeepEnsemble；
    \item \textbf{循环}：PSO基于精炼后的模型继续搜索，如此迭代。
\end{enumerate}

\subsection{浅层一轮后处理精炼}

本课题实现了浅层（Shallow）的一轮后处理精炼，
即执行一次完整的"PSO探索$\to$不确定性反馈$\to$数据增强$\to$模型重训练"循环。

\begin{table}[H]
\centering
\caption{不确定性引导精炼前后的模型性能对比}
\begin{tabular}{lrrrr}
\toprule
指标 & 基线（精炼前） & 精炼后 & 变化量 & 变化率 \\
\midrule
R$^2$ & 0.9995 & 0.9998 & +0.0003 & +0.03\% \\
MSE & 0.3272 & 0.1444 & $-$0.1828 & $-$55.9\% \\
MAE & 0.2712 & 0.1678 & $-$0.1034 & $-$38.1\% \\
RMSE & 0.5720 & 0.3800 & $-$0.1921 & $-$33.6\% \\
\bottomrule
\end{tabular}
\label{tab:al_metrics}
\end{table}

表\ref{tab:al_metrics}显示，
仅一轮不确定性引导精炼后观察到显著的性能提升：
需要指出，当前实验未设置随机重采样（Random Resampling）的消融对照条件，
因此不能确定性能提升中不确定性引导策略的独特贡献——
任何形式的训练集扩充（包括随机复制样本）都可能带来一定程度的精度改善。
严格的因果归因需要消融对照实验。
MSE从0.3272降至0.1444（降低55.9\%），
MAE降低38.1\%，RMSE降低33.6\%。
这表明PSO探索确实识别了代理模型在OOD和低密度区域的薄弱环节。

需要指出，本精炼方法的性能提升主要来源于训练样本的重加权（Reweighting）
效应——高不确定性区域附近的样本在精炼后被赋予了更高的训练覆盖频率，
使得模型在该区域的拟合更充分，而非引入了数据集之外的新物理信息。
因此55.9\%的MSE降幅应理解为模型在OOD区域拟合质量的改善，
不应被解读为代理模型泛化能力的等量提升。

从方法学角度看，本方法本质属于基于不确定性的局部重采样
（Uncertainty-Guided Local Resampling）策略——
其数学效果与对训练样本施加基于$\sigma$的实例权重（Instance Weighting）接近。
采用复制样本而非显式权重的原因在于DeepEnsemble的训练管线
以样本为基本单位，直接赋权需修改损失函数的计算语义。
需要进一步指出，如果检索到的最近邻样本原本已存在于训练集中，
那么本方法的效果等价于对这些样本施加更高的训练频率（即隐式样本重加权），
而非向模型注入了新的数据分布信息。
这一局限性进一步支持了将本方法定性为
局部重采样策略而非信息增强策略的定位。

\subsection{候选解特征分析}

不确定性引导精炼选取了20个高不确定性候选解，
选取100个最近邻训练样本进行数据增强。
Top-5高不确定性候选解的特征分布如下：

\begin{table}[H]
\centering
\caption{Top-5高不确定性候选解}
\begin{tabular}{rrrrrrr}
\toprule
速度 (km/h) & 翼角 (°) & DRS & 下压力 (N) & 阻力 (N) & $\mu$ & $\sigma$ \\
\midrule
175 & 5.1 & 0 & 8773 & 154.8 & 36.34 & 37.57 \\
217 & 14.7 & 1 & 6081 & 260.0 & 46.02 & 36.73 \\
139 & 34.6 & 0 & 6911 & 196.6 & 41.63 & 36.63 \\
106 & 40.4 & 0 & 6331 & 203.4 & 36.17 & 32.13 \\
114 & 22.3 & 0 & 9093 & 155.3 & 27.68 & 31.36 \\
\bottomrule
\end{tabular}
\end{table}

这些高不确定性候选解的共同特征是：
下压力极高（6000--9000N），
对应着训练数据中极少出现的极端工况。
该精炼方法通过在训练集中补充这些区域的相似样本，
有效降低了代理模型在极端工况下的预测不确定性。

\subsection{PSO收敛行为改善}

\begin{figure}[H]
\centering
\includegraphics[width=0.78\textwidth]{figures/active_learning_convergence.png}
\caption{不确定性引导精炼前后的PSO收敛曲线对比}
\label{fig:al_conv}
\end{figure}

图\ref{fig:al_conv}对比了使用基线模型和精炼后的PSO收敛行为。
精炼后的模型使PSO收敛速度提升2.1倍（37 $\to$ 18代），
且最终gbest从101.30（可能含幻觉）降至100.02（更可信）。
这说明：

\begin{enumerate}
    \item 精炼后模型在OOD区域的不确定性降低，
        粒子无需因$\sigma(x)$惩罚而"绕行"，
        直接沿更陡峭的适应度梯度收敛；
    \item 收敛后gbest的降低（101.30 $\to$ 100.02）并非性能退步，
        而是幻觉纠正——基线模型在OOD区域输出了不可靠的虚高评分。
\end{enumerate}

\subsection{小结}

不确定性引导的模型精炼验证了"PSO探索$\to$不确定性反馈$\to$数据增强$\to$模型改进"
单轮反馈的有效性。
需要指出，本实验仅执行了一轮"PSO→增强→重训练"的反馈，
尚未形成多轮迭代的完整闭环。
因此严格来说，已验证的是单轮反馈策略的有效性，
闭环机制的完整验证（多轮迭代收敛性分析及边际收益递减特性）
留给未来工作。
一轮精炼即可将MSE降低55.9\%、收敛速度提升2.1倍，
展示了该机制在降低代理模型OOD不确定性方面的显著价值。
